\documentclass[11pt]{article}

\usepackage[utf8]{inputenc}
\usepackage[T1]{fontenc}
\usepackage[ngerman]{babel}
\usepackage{amsmath}
\usepackage[a4paper,margin=2.5cm]{geometry}
\usepackage{enumitem}
\usepackage[hidelinks]{hyperref}

\newcommand{\paperlabel}[1]{\texttt{#1}}
\newcommand{\papercite}[1]{\texttt{\textbackslash cite\{#1\}}}

\title{Codex-Review der Opus-4.8-Änderungen\\
\large Hinweise zu \texttt{rational\_cut\_and\_project\_gap\_periods.tex}}
\author{}
\date{30. Mai 2026}

\begin{document}
\maketitle

\section*{Drei konkrete Findings}

\begin{enumerate}[leftmargin=2em]
  \item \textbf{Realisierung jedes \(r\in\mathcal R\) etwas sauberer begründen.}

        In der Passage um Zeile~670 ist die neue Aussage mathematisch richtig,
        aber die Formulierung
        \[
          \text{``the range computation of Step~1 shows the strip condition is
          satisfiable for every such \(r\)''}
        \]
        begründet nicht ganz das, was sie begründen soll.  Dass jedes
        \(r\in\mathcal R\) tatsächlich zu akzeptierten Gitterpunkten führt,
        folgt aus der zuvor konstruierten Kongruenzklasse \(P_r\), der
        Integrality-Rechnung über \(\gcd(\alpha,\beta)=1\) bzw. die
        Invertierbarkeit modulo \(D\), und erst danach aus \(r\in\mathcal R\)
        für die Strip-Bedingung.

        Bessere Ersatzformulierung:
        \begin{quote}
        Conversely, if \(\tilde p\in P_r\), then the congruence
        \paperlabel{eq:residue} makes both coordinates in
        \paperlabel{eq:solution}
        integral.  Since \(r\in\mathcal R\), the resulting lattice point
        satisfies the strip condition \paperlabel{eq:r\_range}.  Thus every
        \(r\in\mathcal R\) is realized by accepted lattice points, and
        \(P_r\) is exactly the set of \(\tilde p\)-values with invariant \(r\).
        \end{quote}

  \item \textbf{Randkonvention nicht als exakt ``one less'' formulieren.}

        In der Passage um Zeile~213 ist die neue Erklärung zur geschlossenen
        Intervallkonvention hilfreich.  Die Wendung
        \[
          \text{``rather than one less''}
        \]
        ist aber etwas zu stark: Bei halb-offenen oder anders verschobenen
        Randkonventionen ist die Zählung nicht in jeder Situation schlicht
        exakt um eins kleiner.  Gemeint ist nur, dass die gewählte geschlossene
        Konvention den \(+1\)-Term in
        \[
          N=\lfloor\omega\alpha\rfloor+\lfloor\omega\beta\rfloor+1
        \]
        erklärt.

        Bessere Ersatzformulierung:
        \begin{quote}
        The interval is closed: lattice points lying on either bounding line
        are accepted.  This boundary convention is reflected in the \(+1\) term
        in
        \(N=\lfloor\omega\alpha\rfloor+\lfloor\omega\beta\rfloor+1\), the count
        of admissible levels used in Section~\paperlabel{sec:mainresult}.
        \end{quote}

  \item \textbf{Related-Work-Satz zu Haynes et al. typographisch glätten.}

        Der neue Hinweis auf Haynes--Koivusalo--Walton--Sadun ist inhaltlich
        gut und sollte bleiben.  Um Zeile~1905 entsteht aber eine sehr lange
        Satzfolge:
        \[
          \text{``\dots obtained here is a different invariant. Where the rational case \dots''}
        \]
        Das kompiliert problemlos, liest sich aber im Quelltext und im
        Abschnitt etwas gedrängt.

        Bessere Ersatzformulierung:
        \begin{quote}
        The closest work in the gap-statistics direction is
        \papercite{HaynesEtAl2016}, which studies gap and patch-frequency
        distributions for irrational cut-and-project sets.  The rational,
        sorted-gap minimal period obtained here is a different invariant.

        Where the rational case is treated, the emphasis is on different
        invariants: for instance, \papercite{GaehlerHuntonKellendonk2013}
        computes the integral \v{C}ech cohomology of the tiling spaces
        associated with rational projection-method patterns, an
        algebraic-topological invariant unrelated to the minimal period of the
        sorted gap sequence studied here.
        \end{quote}
\end{enumerate}

\section*{Kurzbilanz}

Keine der drei Stellen ist ein harter mathematischer Fehler.  Die Änderungen
von Opus~4.8 sind insgesamt plausibel; die obigen Punkte sind Präzisierungen,
die den Text robuster und reviewer-fester machen.

\end{document}
